%!TEX root = ../template.tex
%%%%%%%%%%%%%%%%%%%%%%%%%%%%%%%%%%%%%%%%%%%%%%%%%%%%%%%%%%%%%%%%%%%%
%% chapter1.tex
%% NOVA thesis document file
%%
%% Chapter 1: Introduction
%%%%%%%%%%%%%%%%%%%%%%%%%%%%%%%%%%%%%%%%%%%%%%%%%%%%%%%%%%%%%%%%%%%%

\typeout{NT FILE chapter1.tex}%

\chapter{Introduction}
\label{cha:introduction}

\section{Motivation}
\label{sec:intro_motivation}

Software systems are complex socio-technical artefacts that must evolve under technical, organisational, and human constraints. In this setting, software and conceptual modelling provide structured abstractions that can support reasoning, communication, and early validation before implementation. In higher education, modelling is therefore commonly taught in software engineering and related programmes through notations and methods such as UML, BPMN, ER, or goal-oriented approaches.

Despite its recognised importance, modelling is frequently perceived by students in polarised ways: as a valuable instrument for understanding and communicating about systems, or as burdensome ``extra documentation'' that competes with implementation time. These perceptions are not only individual preferences; they can be shaped by socio-technical conditions such as the modelling tools used, the authenticity and domain of assignments, the form and timing of feedback, and the roles that models play in a course (e.g., communication versus strict code generation). Understanding how such factors interact with conceptual and procedural aspects of modelling is important for improving modelling education and for fostering modelling competence that transfers beyond individual courses.

\section{Problem statement and research gap}
\label{sec:intro_gap}

Research on modelling education has produced valuable insights, but it is dispersed across modelling families and research communities. Landscape studies highlight fragmentation in terminology, theoretical grounding, and evaluation approaches across subdomains. In parallel, Bloom-based frameworks and modelling-specific taxonomies provide systematic ways to describe modelling learning outcomes and competencies, often indicating imbalances such as underrepresentation of evaluation and metacognitive aspects. Empirical studies of particular modelling methods further show that beginners face challenges spanning construct semantics, application to requirements, and model quality judgement, and that these challenges can be linked to specific pedagogical and tool-supported requirements.

However, there remains a need to consolidate these perspectives into an integrated understanding of modelling education challenges and support needs that reflects both student and instructor viewpoints across educational contexts.

\section{Aim and research questions}
\label{sec:intro_aim_rq}

The aim of this thesis is to develop an empirically grounded account of modelling education challenges and support needs in higher education, treating modelling as a socio-technical activity that combines conceptual understanding, procedural strategies, and contextual influences.

The study is guided by the following research questions:

\begin{itemize}
    \item \textbf{RQ1:} What challenges do students and instructors perceive in learning and teaching software modelling in higher education, and how do these challenges manifest across different modelling activities (e.g., interpreting, creating, refining, and evaluating models)?
    \item \textbf{RQ2:} What strategies do students and instructors employ to address these challenges, and what support needs (pedagogical and tool-related) do they identify?
    \item \textbf{RQ3:} How do socio-technical factors (e.g., tools, course design, domains, feedback practices) shape modelling learning and teaching experiences?
\end{itemize}

\section{Research approach overview}
\label{sec:intro_approach}

To address these questions, the thesis adopts a qualitative, theory-building approach based on Grounded Theory, informed by Socio-Technical Grounded Theory to explicitly account for interactions between social and technical elements. Data are collected through semi-structured interviews with students and instructors complemented by activity-centered elicitation anchored in concrete modelling tasks or artefacts. This design is intended to elicit both reflective accounts and situated examples of modelling practice and difficulties. Bloom's revised taxonomy is used as a sensitising concept to support interpretation and communication of the cognitive nature of challenges, without constraining the grounded analysis.

\section{Expected contributions}
\label{sec:intro_contributions}

%TODO: Adjust contributions based on actual results


\section{Thesis structure}
\label{sec:intro_structure}

The remainder of this thesis is structured as follows. Chapter~\ref{cha:background} provides background on software and conceptual modelling, modelling education, grounded theory, Bloom-based learning outcomes, and human factors relevant to modelling learning. Chapter~\ref{cha:related_work} reviews related work on modelling education frameworks, empirical studies of modelling learning challenges, tooling and model-driven engineering in education, and contextual factors such as motivation, inclusion, and authenticity. Chapter~\ref{cha:methodology} presents the research design, research questions, data collection and analysis procedures, and considerations of rigour, ethics, and limitations. Subsequent chapters present the findings and discuss implications for modelling education and tool support.
